%% -*- latex-command: lualatex -*-
%
% SPDX-License-Identifier: CC-BY-SA-4.0
% SPDX-FileCopyrightText: Louis Royer <infos.louis.royer@gmail.com>
%
%
% To compile locally (outside of ShareLaTeX), use the following command:
%   make
%
% If compilation fails or if you have rendering issues,
% make sure to have the following packages installed (debian 12):
%   make biber texlive-lualatex texlive-bibtex-extra texlive-lang-french texlive-fonts-extra
%
% To compile inside ShareLaTeX, make sure to set the latex engine to "lualatex".
% This is the most modern version of latex engine currently, and it allows to create lua packages,
% which have many benefits in comparison with LaTeX packages (very hard to write, read, and debug)
% or call to python script (need to pass `--shell-escape` argument to latex engine, and very slow compilation
% (so slow it can fails to build in ShareLaTeX when the document is big,
% even when timeout has been disabled in ShareLaTeX),
% no need to use external files to keep track of state between call to commands)

% Chapters can be individually compiled (to make editing easier), but:
% - links to glossary will be disabled (because glossary chapter will not be included) 
% - first usage of glossary term within the chapter will be full expanded
%   (while with a full compilation only the first usage of within the full document will be expanded)
%   resulting in some page layout alteration
% - page headers/numbering will not be the same as with the full compilation
%   (so you should **not** compile each chapter individually and concatenate them).
% - bibliography numbering will not be the same as wih the full compilation
%
% When creating new chapters, you must make the basename of each chapter different
% (you cannot name each file `main.tex`, and have only the directory being different). 
% Failure to do that will almost certainly result in errors.
% Also, with the current version of luaglossary, there is some limitations on path depth:
% - at root of the project: OK
% - inside a folder: OK (but you should call `\initSufile` at start of document)
% - inside a subfolder: NOT OK

\documentclass{These}

% Bibliography
\addbibresource{biblio.bib}

\begin{document}

\title{Architecture de contrôle pour l’intégration de l’\textit{edge computing} dans les réseaux mobiles 5G}
\author{Louis Royer}

\makeatletter
\hypersetup{
    pdftitle={\@title},
    pdfauthor={\@author},
    pdfsubject={Informatique et télécommunications},
    pdfkeywords={5G; Multi-Access Edge Computing; réseaux mobiles; Mobiles Networks ; redirection de trafic ; Traffic Steering ; gestion de ressources ; Resource Management ; routage par segments ; Segment Routing},
    bookmarksnumbered=true
}
\pdfbookmark{\@title}{title}
\makeatother

\couverture{00_preface/couverture_these.pdf} % page de garde à générer sur ADUM

\frontmatter % page number style = roman

% abstract
\subfile{00_preface/en_abstract}
\subfile{00_preface/fr_abstract}

% remerciements
\cleardoublepage
\subfile{00_preface/remerciements}

% citation
\cleardoublepage
\subfile{00_preface/citation}

% licence
\cleardoublepage
\subfile{00_preface/licence}

% table des matière
\dominitoc
\cleardoublepage % table des matière commence sur une page de droite
\phantomsection
\label{ch:toc}
\pdfbookmark[0]{Table des matières}{ch:toc} % ajout du lien vers la table des matières dans les métadonnées du PDF, mais pas dans le toc imprimé
\tableofcontents

\ContentTitleFormat % mise en forme du titre des chapitres 

\begin{refsection}
% glossaire
\subfile{00_preface/glossary}
\mainmatter % page number style = arabic

% chapitres
\subfile{01_intro/intro}
\subfile{11_overview_5G_edge/overview_5G_edge}
\addtocontents{toc}{\protect\pagebreak} % pour éviter que le titre du 2ème chapitre se retrouve séparé du reste dans le TOC
\subfile{12_state_of_the_art/state_of_the_art}
\addtocontents{toc}{\protect\enlargethispage{2.5\baselineskip}} % pour éviter que la dernière section se retrouve séparée du reste dans le TOC
\addtocontents{toc}{\protect\pagebreak} % pour éviter que le titre du 2ème chapitre se retrouve séparé du reste dans le TOC
\subfile{13_edge_services_integration/edge_services_integration}
\addtocontents{toc}{\protect\vfill} % toc des chapitres 3 et 4 sur la même page, il faut éviter d’avoir des espaces énormes entre les section 4.3.1.1 et 4.3.1.2
\subfile{14_mobility/mobility}
\addtocontents{toc}{\protect\enlargethispage{1\baselineskip}} % pour éviter que la dernière sous-section se retrouve séparée du reste dans le TOC
\subfile{15_testbed_and_eval/testbed_and_eval}
\subfile{20_conclusion/conclusion}

\subfile{30_postface/contributions}

\checkglossaryuse

% bibliographie
\begin{sloppypar} % pour éviter que les URL débordent
\printbibliography
\end{sloppypar}
\addcontentsline{toc}{chapter}{Bibliographie}
\addtocontents{toc}{\protect\clearpage}
\end{refsection}

\PostContentTitleFormat
\quatriemeDeCouverture{30_postface/4e_de_couverture_these.pdf} % page de garde à générer sur ADUM
\end{document}
